% Load-strength interference: two Gaussian densities, load L (left)
% and strength S (right), with the overlap region carrying the
% failure probability. Drawn in scaled stress units; the load is
% centred at 4.0 with sigma 0.6 and the strength at 6.0 with sigma
% 0.8 (arbitrary plot units, matching the relative spread of the
% worked vessel example).
\begin{figure}[htbp]
\centering
\begin{tikzpicture}[x=1.4cm, y=3.0cm, font=\small]

% Axes
\draw[->, thick] (1.4, 0) -- (8.6, 0) node[right] {stress};
\draw[->, thick] (1.4, -0.03) -- (1.4, 1.15) node[above] {density};

% Shaded interference region (overlap = area under the lower envelope,
% the min of the two densities). The two Gaussians cross at x ~ 4.925.
% Left of the crossing the strength density is lower; right of it the
% load density is lower. The shaded mass is the genuine overlap.
\fill[engamber, opacity=0.40]
  plot[domain=4.10:4.925, samples=40]
    (\x, {1.5*0.49868*exp(-0.5*pow((\x-6.0)/0.8,2))})
  -- plot[domain=4.925:6.0, samples=40]
    (\x, {1.5*0.66467*exp(-0.5*pow((\x-4.0)/0.6,2))})
  -- (6.0, 0) -- (4.10, 0) -- cycle;

% Load density: N(4.0, 0.6^2), peak height 1/(0.6 sqrt(2pi)) ~ 0.665,
% rescaled to plot height ~1.0 by factor 1.5.
\draw[thick, engblue, smooth]
  plot[domain=2.0:6.2, samples=100]
  (\x, {1.5*0.66467*exp(-0.5*pow((\x-4.0)/0.6,2))});
\node[anchor=east, engblue] at (3.4, 0.72) {load $L$};

% Strength density: N(6.0, 0.8^2), peak ~ 0.4986, rescaled by 1.5.
\draw[thick, engred, smooth]
  plot[domain=3.6:8.4, samples=100]
  (\x, {1.5*0.49868*exp(-0.5*pow((\x-6.0)/0.8,2))});
\node[anchor=west, engred] at (6.8, 0.62) {strength $S$};

% Interference annotation
\node[anchor=west, engamber!80!black, font=\scriptsize]
  at (6.6, 0.32) {interference $\Rightarrow p_{f}$};
\draw[->, engamber!80!black]
  (6.6, 0.30) .. controls (5.9, 0.20) .. (5.05, 0.13);

% Mean markers
\draw[dashed, engblue!70] (4.0, 0) -- (4.0, 0.92);
\draw[dashed, engred!70] (6.0, 0) -- (6.0, 0.69);
\node[anchor=north, font=\scriptsize] at (4.0, -0.02) {$\mu_L$};
\node[anchor=north, font=\scriptsize] at (6.0, -0.02) {$\mu_S$};

% margin annotation
\draw[<->, thick, engblack!70] (4.0, 1.04) -- (6.0, 1.04);
\node[anchor=south, font=\scriptsize] at (5.0, 1.04)
  {margin $\mu_S - \mu_L$};

\end{tikzpicture}
\caption[Load-strength interference]{The load-strength interference
model. The component carries a random load $L$ and has a random
strength $S$, both in the same stress units. Failure is the event
$L \geq S$, and its probability is the mass in the shaded overlap of
the two densities, not the gap between the two means. Two designs with
the same nominal safety factor $\mu_S/\mu_L$ can have failure
probabilities orders of magnitude apart, because the overlap is set by
the spreads $\sigma_S$ and $\sigma_L$ as much as by the means. For
independent Gaussian inputs the margin $M = S - L$ is Gaussian and the
failure probability is $\Phi(-\beta)$ with reliability index $\beta =
(\mu_S - \mu_L)/\sqrt{\sigma_S^2 + \sigma_L^2}$. Source: schematic by
the authors; worked numbers in the text.}
\label{fig:vol02:ch13:load-strength}
\end{figure}
